\chapter{Introduction}

Graph-structured data appears throughout machine learning, including citation networks,
social networks, biological interaction networks, and knowledge graphs. In these settings,
the predictive signal is not contained only in node features. It is also carried by the
edges that determine which nodes exchange information. Graph Neural Networks (GNNs) exploit
this relational structure through message passing, where each node updates its representation
by aggregating information from its neighbours \cite{kipf2017gcn, hamilton2017graphsage}.
This has made GNNs a standard approach for semi-supervised node classification.

The benefit of message passing depends strongly on graph quality. If edges connect nodes
with related labels or compatible features, aggregation can reinforce useful information.
If the graph contains noisy, missing, or misleading edges, the same aggregation mechanism can
propagate the wrong signal. This makes graph topology an important part of the learning
pipeline rather than a fixed background object. Improving the graph can therefore be viewed
as a complementary strategy to improving the GNN architecture itself.

One structural property closely tied to GNN performance is \textit{homophily}, the tendency
of connected nodes to share the same class. Standard message-passing GNNs often perform well
on homophilous graphs because neighbours provide class-consistent information. On graphs
with weak homophily, aggregation can mix incompatible signals and reduce performance
\cite{zhu2020homophily}. Despite this relationship, many graph structure learning and graph
editing methods do not treat homophily as an explicit reward or objective. Homophily is often
left as an implicit side effect of optimizing a task loss or reconstruction objective.

This thesis studies whether homophily can be used directly as a reward signal for graph
editing. We propose \textbf{GraphHARE} (Graph Homophily-Aware Reinforcement Editing), a
reinforcement learning framework that edits the topology of a citation graph while keeping a
pre-trained GraphSAGE classifier fixed. The agent receives feedback from a held-out validation
split and from the structural change in graph homophily. The intended reward has the form
\begin{equation}
    r_t =
    \frac{\Delta F_1}{s_{F_1}}
    + \beta \frac{\Delta h}{s_h},
    \label{eq:reward}
\end{equation}
where $\Delta F_1$ is the change in validation macro-F1 between consecutive evaluation
points, $\Delta h$ is the change in reward-safe homophily, $s_{F_1}$ and $s_h$ are fixed
calibration scales, and $\beta \geq 0$ controls the weight of the homophily term. The frozen
backbone design separates graph editing from classifier training, making the experiment a
test of whether topology edits alone can improve a fixed GNN.

A major part of this work became methodological. Early single-seed results appeared to show
that GraphHARE improved Cora and that performance followed an inverted-U curve over $\beta$.
A later audit found that those results were not valid. Two issues were responsible. First,
the homophily reward used ground-truth labels on test-incident edges, leaking information
from the test set into the reward for every $\beta > 0$ run. Second, the reported headline
metric was the maximum validation macro-F1 over hundreds of noisy episodes, which selected
the luckiest validation fluctuation rather than a stable learned improvement. After fixing
the reward, evaluating test performance once after training, and reporting 10-seed paired
confidence intervals, the apparent gains disappeared.

The corrected experiments on Cora, CiteSeer, and PubMed show that GraphHARE matches the
frozen GraphSAGE baseline, with no statistically significant test improvement for
$\beta \in \{0, 0.5, 1.0\}$. This null result is important because the baselines themselves
remain sound: GraphSAGE strongly outperforms an MLP on Cora, moderately outperforms it on
CiteSeer, and only slightly outperforms it on PubMed. The graph is useful to the classifier,
but the tested reward and policy do not find additional test-transferable improvements on
these already homophilous citation graphs.

To understand this outcome, we also perform an oracle headroom analysis. The analysis asks
how much test performance could improve if graph edits were chosen with different amounts of
information. A full oracle using true labels can produce large gains, showing that topology
can matter in principle. However, when the editor is restricted to information available to
the actual policy, the reachable gain is small on Cora and approximately zero on CiteSeer
and PubMed. This suggests that the failure is not only an optimization problem. On these
datasets, most of the useful editing signal is not available without label information that
the protocol correctly forbids.

\section*{Contributions}

The main contributions of this thesis are:

\begin{itemize}
    \item A homophily-aware reinforcement learning formulation for graph editing, combining
    validation macro-F1 and reward-safe homophily in a calibrated scalar reward.

    \item A frozen-backbone graph editing protocol using GraphSAGE, designed to test whether
    topology refinement alone can improve a fixed node classifier.

    \item A corrected evaluation protocol that separates training, model selection, reward
    evaluation, and test reporting, with 10-seed paired confidence intervals.

    \item An empirical null result on Cora, CiteSeer, and PubMed showing that homophily-aware
    RL graph editing does not significantly improve test performance over frozen GraphSAGE
    on these homophilous citation benchmarks.

    \item An oracle headroom analysis explaining the null result by separating theoretical
    topology headroom from prediction-reachable editing headroom.

    \item An expanded comparison harness including MLP, GraphSAGE, GCN, GAT, APPNP, GCNII,
    GraphTrans, and a graph autoencoder baseline.
\end{itemize}

\section*{Thesis Outline}

Chapter~2 reviews graph neural networks, homophily, graph autoencoders, reinforcement
learning for graph problems, and graph topology optimisation. Chapter~3 defines the
GraphHARE environment, reward function, policy network, and corrected evaluation protocol.
Chapter~4 presents the experimental results, baseline comparisons, and oracle headroom
analysis. Chapter~5 concludes with lessons from the null result and directions for future
work on graph editing.
